% ============================================================
%  chapter/chapter1.tex — Mathematics stress test
% ============================================================
\chapter{Mathematics}\label{ch:maths}

Inline maths such as $e^{i\pi} + 1 = 0$ sits in the running text;
display maths gets its own line and number:

\begin{equation}\label{eq:gauss}
  \int_{-\infty}^{\infty} e^{-x^{2}} \,\mathrm{d}x = \sqrt{\pi}.
\end{equation}

Equation~\eqref{eq:gauss} is referenced with \texttt{\textbackslash eqref}. Multi-line
derivations use \texttt{align}:

\begin{align}
  (a+b)^2 &= a^2 + 2ab + b^2,\\
  (a-b)^2 &= a^2 - 2ab + b^2,\\
  (a-b)(a+b) &= a^2 - b^2.
\end{align}

Matrices and cases work through \texttt{amsmath}:

\begin{equation}
  A = \begin{pmatrix}
        a_{11} & a_{12} & \cdots & a_{1n}\\
        a_{21} & a_{22} & \cdots & a_{2n}\\
        \vdots & \vdots & \ddots & \vdots\\
        a_{m1} & a_{m2} & \cdots & a_{mn}
      \end{pmatrix},
  \qquad
  f(x) = \begin{cases}
           x^{2}, & x \geq 0,\\
           -x,    & x < 0.
         \end{cases}
\end{equation}

Theorem-like statements can be built with the classic
\texttt{amsthm} machinery if your book needs it; for most trade
books the styled paragraph below is enough.

\section{Symbols in the index}
Technical terms should reach the index as you introduce them: the
Euler identity\index{Euler identity}, Gaussian
integrals\index{Gaussian integral} and
matrices\index{matrix|see{linear algebra}} all appear above.

\section{Statements and proofs}
Formal books state, then prove. The environments come from
\emph{amsthm}, styled in this template's navy-and-sans scheme:

\begin{definition}\label{def:prime}
A natural number $p > 1$ is \emph{prime} if its only positive
divisors are $1$ and $p$ itself.
\end{definition}

\begin{theorem}[Euclid]\label{thm:primes}
There are infinitely many primes.
\end{theorem}

\begin{proof}
Suppose $p_1, p_2, \dots, p_n$ were all the primes and consider
$N = p_1 p_2 \cdots p_n + 1$. No $p_i$ divides $N$, since each
leaves remainder~$1$; yet $N > 1$ has a prime factor by
Definition~\ref{def:prime}. Contradiction.
\end{proof}

\begin{example}
The first few primes are $2, 3, 5, 7, 11, 13$;
by Theorem~\ref{thm:primes} the list never ends.\index{prime number}
\end{example}

\begin{note}
Unnumbered remarks use the starred \texttt{note} environment ---
useful for asides that would clutter the numbering.
\end{note}

\section{More display mathematics}
Fractions, binomials and named operators behave as usual:
\begin{equation}\label{eq:binom}
  \binom{n}{k} = \frac{n!}{k!\,(n-k)!},
  \qquad \sum_{k=0}^{n} \binom{n}{k} = 2^{n},
  \qquad \lim_{x \to 0} \frac{\sin x}{x} = 1.
\end{equation}
Text belongs inside maths through \texttt{\textbackslash text},
which keeps the surrounding font:
\begin{equation*}
  f(x) = x^{2} \quad \text{for all } x \in \mathbb{R}.
\end{equation*}
Both numbered forms cross-reference: combine
\eqref{eq:gauss} with \eqref{eq:binom} and you have most of a
first calculus course.\index{binomial coefficient}
